
\documentclass{article}
\usepackage{colortbl}
\usepackage{makecell}
\usepackage{multirow}
\usepackage{supertabular}

\begin{document}

\newcounter{utterance}

\centering \large Interaction Transcript for game `dond', experiment `semi\_it', episode 2 with gpt{-}5{-}mini{-}2025{-}08{-}07{-}no{-}reasoning{-}t1.0.
\vspace{24pt}

{ \footnotesize  \setcounter{utterance}{1}
\setlength{\tabcolsep}{0pt}
\begin{supertabular}{c@{$\;$}|p{.15\linewidth}@{}p{.15\linewidth}p{.15\linewidth}p{.15\linewidth}p{.15\linewidth}p{.15\linewidth}}
    \# & $\;$A & \multicolumn{4}{c}{Game Master} & $\;\:$B\\
    \hline

    \theutterance \stepcounter{utterance}  
    & & \multicolumn{4}{p{0.6\linewidth}}{
        \cellcolor[rgb]{0.9,0.9,0.9}{
            \makecell[{{p{\linewidth}}}]{
                \texttt{\tiny{[P1$\langle$GM]}}
                \texttt{State giocando a un gioco di negoziazione in cui dovete accordarvi su come dividere una serie di oggetti tra voi e un altro giocatore.} \\
\\ 
\texttt{Regole:} \\
\texttt{(a) A Lei e all'altro giocatore viene dato un insieme di oggetti. Ognuno di voi riceve anche una funzione di valore segreta, che rappresenta il valore di ciascun tipo di oggetto.} \\
\texttt{(b) Si scambiano messaggi con l'altro giocatore per concordare chi si aggiudica gli oggetti. Potete inviare un massimo di 5 messaggi ciascuno, oppure terminare in anticipo facendo una proposta segreta in qualsiasi momento.} \\
\texttt{(c) A ciascuno di voi viene chiesto di inviare una proposta segreta indicando gli oggetti che desiderate, formattata tra parentesi quadre come segue: "{[}Proposta: <numero> <nome oggetto>, <numero> <nome oggetto>, <...>{]}".} \\
\texttt{(d) Se le vostre proposte sono complementari, cioè ci sono abbastanza oggetti per soddisfare entrambe le proposte, a ciascun giocatore viene assegnato un punteggio basato sulla somma dei valori degli oggetti ricevuti. In caso contrario, entrambi ricevono zero punti.} \\
\texttt{(e) L'obiettivo è massimizzare il proprio punteggio.} \\
\\ 
\texttt{Cominciamo.} \\
\\ 
\texttt{L'insieme degli oggetti disponibili è:} \\
\\ 
\texttt{1 pala, 2 bacche, 1 ciotola, 3 bicchieri} \\
\\ 
\texttt{La funzione valore segreta è:} \\
\\ 
\texttt{pala: 6, bacca: 0, ciotola: 1, bicchiere: 1} \\
\\ 
\texttt{IMPORTANTE: i vostri messaggi, a meno che non si tratti di una proposta segreta, vengono trasmessi direttamente all'altro giocatore, quindi non includete alcuna risposta alle regole o testo di annuncio del vostro messaggio. Per fare una proposta segreta, utilizzate il formato indicato. Non utilizzare delle parentesi quadre quando si comunica all'altro giocatore, altrimenti verrà interpretata come una proposta segreta.} \\
            }
        }
    }
    & \\ \\

    \theutterance \stepcounter{utterance}  
    & \multicolumn{4}{p{0.6\linewidth}}{
        \cellcolor[rgb]{0.9,0.9,0.9}{
            \makecell[{{p{\linewidth}}}]{
                \texttt{\tiny{[P1$\rangle$GM]}}
                \texttt{Ciao! Per me la pala vale molto (6), i bicchieri e la ciotola valgono poco (1 ciascuno), le bacche zero. Ti proporrei: io prendo la pala, e tu prendi le 2 bacche. Possiamo poi dividere ciotola e bicchieri in base a cosa ti serve di più. Cosa ti interessa di più tra ciotola e bicchieri?} \\
            }
        }
    }
    & & \\ \\

    \theutterance \stepcounter{utterance}  
    & & & \multicolumn{2}{p{0.3\linewidth}}{
        \cellcolor[rgb]{0.9,0.9,0.9}{
            \makecell[{{p{\linewidth}}}]{
                \texttt{\tiny{[GM$|$GM]}}
                \texttt{continue} \\
            }
        }
    }
    & & \\ \\

    \theutterance \stepcounter{utterance}  
    & & & \multicolumn{2}{p{0.3\linewidth}}{
        \cellcolor[rgb]{0.9,0.9,0.9}{
            \makecell[{{p{\linewidth}}}]{
                \texttt{\tiny{[GM$|$GM]}}
                \texttt{Ciao! Per me la pala vale molto (6), i bicchieri e la ciotola valgono poco (1 ciascuno), le bacche zero. Ti proporrei: io prendo la pala, e tu prendi le 2 bacche. Possiamo poi dividere ciotola e bicchieri in base a cosa ti serve di più. Cosa ti interessa di più tra ciotola e bicchieri?} \\
            }
        }
    }
    & & \\ \\

    \theutterance \stepcounter{utterance}  
    & & \multicolumn{4}{p{0.6\linewidth}}{
        \cellcolor[rgb]{0.9,0.9,0.9}{
            \makecell[{{p{\linewidth}}}]{
                \texttt{\tiny{[GM$\rangle$P2]}}
                \texttt{State giocando a un gioco di negoziazione in cui dovete accordarvi su come dividere una serie di oggetti tra voi e un altro giocatore.} \\
\\ 
\texttt{Regole:} \\
\texttt{(a) A Lei e all'altro giocatore viene dato un insieme di oggetti. Ognuno di voi riceve anche una funzione di valore segreta, che rappresenta il valore di ciascun tipo di oggetto.} \\
\texttt{(b) Si scambiano messaggi con l'altro giocatore per concordare chi si aggiudica gli oggetti. Potete inviare un massimo di 5 messaggi ciascuno, oppure terminare in anticipo facendo una proposta segreta in qualsiasi momento.} \\
\texttt{(c) A ciascuno di voi viene chiesto di inviare una proposta segreta indicando gli oggetti che desiderate, formattata tra parentesi quadre come segue: "{[}Proposta: <numero> <nome oggetto>, <numero> <nome oggetto>, <...>{]}".} \\
\texttt{(d) Se le vostre proposte sono complementari, cioè ci sono abbastanza oggetti per soddisfare entrambe le proposte, a ciascun giocatore viene assegnato un punteggio basato sulla somma dei valori degli oggetti ricevuti. In caso contrario, entrambi ricevono zero punti.} \\
\texttt{(e) L'obiettivo è massimizzare il proprio punteggio.} \\
\\ 
\texttt{Cominciamo.} \\
\\ 
\texttt{L'insieme degli oggetti disponibili è:} \\
\\ 
\texttt{1 pala, 2 bacche, 1 ciotola, 3 bicchieri} \\
\\ 
\texttt{La funzione valore segreta è:} \\
\\ 
\texttt{pala: 4, bacca: 1, ciotola: 4, bicchiere: 0} \\
\\ 
\texttt{IMPORTANTE: i vostri messaggi, a meno che non si tratti di una proposta segreta, vengono trasmessi direttamente all'altro giocatore, quindi non includete alcuna risposta alle regole o testo di annuncio del vostro messaggio. Per fare una proposta segreta, utilizzate il formato indicato. Non utilizzare delle parentesi quadre quando si comunica all'altro giocatore, altrimenti verrà interpretata come una proposta segreta.} \\
\\ 
\\ 
\texttt{Ciao! Per me la pala vale molto (6), i bicchieri e la ciotola valgono poco (1 ciascuno), le bacche zero. Ti proporrei: io prendo la pala, e tu prendi le 2 bacche. Possiamo poi dividere ciotola e bicchieri in base a cosa ti serve di più. Cosa ti interessa di più tra ciotola e bicchieri?} \\
            }
        }
    }
    & \\ \\

    \theutterance \stepcounter{utterance}  
    & & & \multicolumn{4}{p{0.6\linewidth}}{
        \cellcolor[rgb]{0.9,0.9,0.9}{
            \makecell[{{p{\linewidth}}}]{
                \texttt{\tiny{[GM$\langle$P2]}}
                \texttt{Per me la ciotola vale molto e i bicchieri nulla. Accetto che tu prenda la pala; io vorrei la ciotola. Possiamo lasciarti tutti e 3 i bicchieri se per te hanno valore, e tu puoi prenderti anche le 2 bacche visto che per me valgono poco. Va bene: tu pala, bicchieri e bacche; io ciotola?} \\
            }
        }
    }
     \\ \\

    \theutterance \stepcounter{utterance}  
    & & & \multicolumn{2}{p{0.3\linewidth}}{
        \cellcolor[rgb]{0.9,0.9,0.9}{
            \makecell[{{p{\linewidth}}}]{
                \texttt{\tiny{[GM$|$GM]}}
                \texttt{continue} \\
            }
        }
    }
    & & \\ \\

    \theutterance \stepcounter{utterance}  
    & & & \multicolumn{2}{p{0.3\linewidth}}{
        \cellcolor[rgb]{0.9,0.9,0.9}{
            \makecell[{{p{\linewidth}}}]{
                \texttt{\tiny{[GM$|$GM]}}
                \texttt{Per me la ciotola vale molto e i bicchieri nulla. Accetto che tu prenda la pala; io vorrei la ciotola. Possiamo lasciarti tutti e 3 i bicchieri se per te hanno valore, e tu puoi prenderti anche le 2 bacche visto che per me valgono poco. Va bene: tu pala, bicchieri e bacche; io ciotola?} \\
            }
        }
    }
    & & \\ \\

    \theutterance \stepcounter{utterance}  
    & & \multicolumn{4}{p{0.6\linewidth}}{
        \cellcolor[rgb]{0.9,0.9,0.9}{
            \makecell[{{p{\linewidth}}}]{
                \texttt{\tiny{[P1$\langle$GM]}}
                \texttt{Per me la ciotola vale molto e i bicchieri nulla. Accetto che tu prenda la pala; io vorrei la ciotola. Possiamo lasciarti tutti e 3 i bicchieri se per te hanno valore, e tu puoi prenderti anche le 2 bacche visto che per me valgono poco. Va bene: tu pala, bicchieri e bacche; io ciotola?} \\
            }
        }
    }
    & \\ \\

    \theutterance \stepcounter{utterance}  
    & \multicolumn{4}{p{0.6\linewidth}}{
        \cellcolor[rgb]{0.9,0.9,0.9}{
            \makecell[{{p{\linewidth}}}]{
                \texttt{\tiny{[P1$\rangle$GM]}}
                \texttt{Perfetto, accordo ottimo per me: io prendo pala, 3 bicchieri e 2 bacche; tu prendi la ciotola.} \\
\\ 
\texttt{{[}Proposta: 1 pala, 2 bacche, 3 bicchieri{]}} \\
            }
        }
    }
    & & \\ \\

    \theutterance \stepcounter{utterance}  
    & & & \multicolumn{2}{p{0.3\linewidth}}{
        \cellcolor[rgb]{0.9,0.9,0.9}{
            \makecell[{{p{\linewidth}}}]{
                \texttt{\tiny{[GM$|$GM]}}
                \texttt{proposal} \\
            }
        }
    }
    & & \\ \\

    \theutterance \stepcounter{utterance}  
    & & & \multicolumn{2}{p{0.3\linewidth}}{
        \cellcolor[rgb]{0.9,0.9,0.9}{
            \makecell[{{p{\linewidth}}}]{
                \texttt{\tiny{[GM$|$GM]}}
                [1, 2, 0, 3]
            }
        }
    }
    & & \\ \\

    \theutterance \stepcounter{utterance}  
    & & \multicolumn{4}{p{0.6\linewidth}}{
        \cellcolor[rgb]{0.9,0.9,0.9}{
            \makecell[{{p{\linewidth}}}]{
                \texttt{\tiny{[GM$\rangle$P2]}}
                \texttt{L'altro giocatore ha presentato una proposta segreta. Ora tocca a Lei presentare una sua proposta.} \\
            }
        }
    }
    & \\ \\

    \theutterance \stepcounter{utterance}  
    & & & \multicolumn{4}{p{0.6\linewidth}}{
        \cellcolor[rgb]{0.9,0.9,0.9}{
            \makecell[{{p{\linewidth}}}]{
                \texttt{\tiny{[GM$\langle$P2]}}
                \texttt{{[}Proposta: 1 pala, 2 bacche, 0 ciotola, 3 bicchieri{]}} \\
            }
        }
    }
     \\ \\

    \theutterance \stepcounter{utterance}  
    & & & \multicolumn{2}{p{0.3\linewidth}}{
        \cellcolor[rgb]{0.9,0.9,0.9}{
            \makecell[{{p{\linewidth}}}]{
                \texttt{\tiny{[GM$|$GM]}}
                \texttt{proposal} \\
            }
        }
    }
    & & \\ \\

    \theutterance \stepcounter{utterance}  
    & & & \multicolumn{2}{p{0.3\linewidth}}{
        \cellcolor[rgb]{0.9,0.9,0.9}{
            \makecell[{{p{\linewidth}}}]{
                \texttt{\tiny{[GM$|$GM]}}
                [1, 2, 0, 3]
            }
        }
    }
    & & \\ \\

    \theutterance \stepcounter{utterance}  
    & & & \multicolumn{2}{p{0.3\linewidth}}{
        \cellcolor[rgb]{0.9,0.9,0.9}{
            \makecell[{{p{\linewidth}}}]{
                \texttt{\tiny{[GM$|$GM]}}
                [[1, 2, 0, 3], [1, 2, 0, 3]]
            }
        }
    }
    & & \\ \\

\end{supertabular}
}

\end{document}
